\section{Multi-Criteria Analysis}\label{sec:multi-criteria-analysis}

\subsection{Criteria}\label{subsec:criteria}
A set of criteria was established to evaluate the hybrid frameworks.
The reasoning behind each criterion, as well as their assigned weights, are explained below.

\subsubsection{Communication Protocols}\label{subsubsec:communication-protocols}
The to be build software solution requires two types of communication.

The back-end application will provide an API that follows the REST protocol standards.
Since the front-end application needs to communicate with the back-end to perform CRUD operations, it must be able to construct and send \emph{RESTful API calls}.

Additionally, the software solution will likely include chat functionality.
For real-time communication between users, some sort of socket connection is required.
This means the front-end must also be able to establish and maintain a \emph{WebSocket connection}.

\paragraph{Method of evaluation}
To score this criterion, packages or built-in solutions will be looked up for both RESTful API calls and WebSocket connections.
In case any of these functionalities are not supported, the framework will be disqualified.

\paragraph{Weight}
This criterion is not assigned any weight.

If the framework does not give the option to perform both acts of communication, the functionality of the software solution cannot be realized, meaning the application cannot be built using that framework.

\subsubsection{Community}\label{subsubsec:community}
The larger the community, the more online resources and support will be available for learning and troubleshooting the framework.
To assess the size of the community, several metrics will be considered:

\begin{itemize}
	\item The \emph{number of stars} on the official framework GitHub repository.
	      While not a direct measure of usage, it gives a rough indication on how many developers are interested in or actively using the framework.
	      The more developers using the framework, the more questions that will have been asked and answered online.
	\item The number of \emph{Discord users} in the official or related Discord servers.
	      A larger user count means more potential help when asking questions.
	\item The number of \emph{examples} online.
	      Examples of implementations, tutorials, and blog posts can be very helpful when learning a new framework.
	      If the framework is unopinionated, a big variety of examples can help in deciding how to structure the application.
\end{itemize}

\paragraph{Method of evaluation}
To score this criterion, the stars from the official GitHub repository will be counted, the official or related Discord servers will be joined to view the member count, and a search on GitHub will be performed to find the number of repositories using the framework.

\paragraph{Weight}
This criterion is assigned a weight of \emph{20}, derived from the above sub-criteria, which are \emph{5}, \emph{5}, and \emph{10}.

Both the number of stars and the number of Discord users are not very important, because an application can be built without community support as well.
It would just be more difficult.

The number of available examples holds slightly more value, as it provides real-world insights into project structure and professional best practices.
That said, while they can speed up the learning process, and application can still be built without them.

For these reasons, the community criterion is not assigned a very high weight.

\subsubsection{Developer Experience}\label{subsubsec:developer-experience}
A good developer experience is important for productivity and ease of use.
Even if a framework is powerful and flexible, if it is difficult to use, it will slow down development and may frustrate the developer.
It should be easy to work with, so developers can focus on building the application rather than struggling with the tooling.
The following aspects will be considered to determine good developer experience:

\begin{itemize}
	\item The framework and the programming language it uses should not have a high \emph{learning curve}.
	      If a framework differs significantly from other frameworks, one which the developer is not yet familiar with, building an application within a short time becomes challenging.
	\item There should be built-in support for \emph{debugging} the application during runtime using a Jetbrains IDE.
	      Without proper debugging tools, identifying and fixing issues can be time-consuming and difficult.
	\item High quality \emph{documentation} should be readily available.
	      Poor documentation forces developers to rely on user discussions and examples, which may not always be accurate or up-to-date.
	      Good documentation should include clear explanations, code examples, and guides on best practices.
\end{itemize}

\paragraph{Method of evaluation}
To score this criterion, the existence of a debugging solution as well as the presence of proper documentation will be checked.
To validate that the documentation is proper, the website must have some sort of installation guide, and documentation on specific parts of the framework, such as state management and navigation.

The other part, the learning curve, will be determined subjectively based on the developer`s prior experience with similar frameworks and programming languages.
The proficiency scale is as follows:

\begin{itemize}
	\item \textbf{0.00} - No prior experience with the framework or programming language.
	\item \textbf{0.25} - Limited experience with the framework or programming language.
	      The developer has worked with it on small projects or has only read about it.
	\item \textbf{0.50} - Moderate experience with the framework or programming language.
	      The developer has worked with it on a few projects and has a basic understanding of its features and best practices.
	\item \textbf{0.75} - Good experience with the framework or programming language.
	      The developer has worked with it on several projects and is familiar with its features and best practices.
	\item \textbf{1.00} - Extensive experience with the framework or programming language.
	      The developer has worked with it on many projects and is considered an expert in its use
\end{itemize}

\paragraph{Weight}
This criterion is assigned a weight of \emph{45}, derived from the above sub-criteria, which are \emph{20}, \emph{10}, and \emph{15}.

The learning curve is quite important.
Time is the limiting factor in this project, and both prototypes and the final application need to be built within a few weeks.
If a framework takes too long to learn, it will not be feasible to use it.

Debugging is also important, but even without it, development is still very doable.
Most programming languages have some sort of logging functionality, which can be used to identify issues instead.

Documentation is very important, especially when working with a framework for the first time.
While a large and proficient community can provide answers, it is better to have standardized information readily available rather than having to rely solely on other users.

For these reasons, the developer experience criterion is assigned a relatively high weight.

\subsubsection{Ecosystem}\label{subsubsec:ecosystem}
The ecosystem comprises all available resources to reduce the workload of developers.
The more resources are available, the more developers can focus on implementing business logic rather than reinventing the wheel.
The following aspects will be considered to determine the quality of the ecosystem:

\begin{itemize}
	\item There should be official or well-maintained community \emph{plugins} for Jetbrains IDEs, offering features such as code completion, tooltip documentation, and error checking.
	      These plugins can significantly enhance productivity and reduce the likelihood of errors.
	\item A wide variety of \emph{libraries} should be available.
	      These libraries can provide pre-built functionality, saving developers time and effort.
	      The higher the number of available libraries, the bigger the chance that a library exists for a specific need.
	\item For the user interface, there should be premade \emph{UI components}, which may be built-in or from open source libraries.
	      These components can help developers to quickly create screens and layouts without having to design everything from scratch.
\end{itemize}

\paragraph{Method of evaluation}
To score this criterion, the availability of IDE plugin support will be checked, as well as the number of existing libraries and the existence of one or more UI component libraries.

\paragraph{Weight}
This criterion is assigned a weight of \emph{55}, derived from the above sub-criteria, which are \emph{25}, \emph{15}, and \emph{15}.

The availability of IDE plugins is very important.
Especially when working with an Android emulator, having tools that streamline development makes a big difference.
If these tools are not available, development will be slower and more error-prone.

The number of libraries and UI components are also important, but to a lesser extent.
While they can speed up development, developers can create their own libraries and components if needed.
This however, takes more time, making a rich ecosystem of ready-to-use libraries and components preferable.

For these reasons, the ecosystem criterion is assigned the highest weight.

\subsubsection{Quality Assurance}\label{subsubsec:quality-assurance}
The app should be of good quality when released to the public, ensuring both solid security and proper functionality.
This criterion encapsulates a part of the advanced track chosen for back-end development (\emph{Quality Management and Deployment}).
To help ensure the quality of the application, the following aspects will be considered:

\begin{itemize}
	\item The framework should provide some form of (automated) \emph{testing}.
	      Testing a mobile front-end framework can be challenging, but if this is not done, new updates may introduce bugs or break existing functionality, which could confuse or frustrate the users.
	\item The framework should provide some form of \emph{static code analysis} tooling.
	      This ensures the code is of good quality and keeps consistency between code that is written by different developers.
	\item The framework should receive regular \emph{updates}.
	      If bugs or security flaws are discovered, they should be fixed timely.
	      Without frequent updates, users are left at risk, while developers have no control of critical issues.
\end{itemize}

\paragraph{Method of evaluation}
To score this criterion, the existence of unit and component testing solutions, and linting tools, will be checked.
In addition, it will be verified that the framework has received an update within the last six months.

\paragraph{Weight}
This criterion is assigned a weight of \emph{40}, derived from the above sub-criteria, which are \emph{10}, \emph{5}, and \emph{25}.

Testing the application is important, since it directly affects the quality of the deployed application from a user`s perspective.

Static code analysis is not very important.
While it helps maintain cleaner and more consistent code, it does not directly affect the quality of the deployed application from a user`s perspective.
Since this project is small-scale, the codebase will be manageable without it.

Regular updates are very important, since it can have substantial consequences for the safety of the user`s data if they are long-term security problems.

For these reasons, the quality assurance criterion is assigned a relatively high weight.

\subsection{Frameworks}\label{subsec:frameworks}
The following hybrid frameworks will be evaluated based on the criteria established in \cref{subsec:criteria}.

\subsubsection{Flutter}\label{subsubsec:frameworks-flutter}
Flutter is an open source framework developed by Google for building natively compiled applications for mobile, web, and desktop from a single codebase. \cite{bib:most-popular-mobile-frameworks-2024}
It uses the Dart programming language and provides a rich set of pre-designed UI components and tools for building high-performance applications.

For the retake, Flutter is a mandatory framework to research, but we had initially chosen it because it is one of the most popular frameworks for mobile development \cite{bib:most-popular-mobile-frameworks-2024}, as well as the fact it uses its own programming language, which makes it different from the other frameworks.

\subsubsection{React Native}\label{subsubsec:frameworks-react-native}
React Native is an open source framework developed by Facebook for building cross-platform mobile applications using JavaScript and React.
It allows developers to build mobile apps that can run on both iOS and Android platforms using a single codebase. \cite{bib:most-popular-mobile-frameworks-2024}

Like Flutter, React Native is now a mandatory framework to research for the retake too.
It was initially chosen because it is the most popular frameworks for mobile development \cite{bib:most-popular-mobile-frameworks-2024}, and because it is similar to web development, which we are already familiar with.

\subsubsection{Ionic}\label{subsubsec:frameworks-ionic}
Ionic is a popular open-source framework for building cross-platform mobile applications using web technologies such as HTML, CSS, and JavaScript.
It provides a library of pre-designed UI components and tools for building mobile apps that can run on both iOS and Android platforms. \cite{bib:most-popular-mobile-frameworks-2024}

Ionic was chosen because it is the most popular PWA framework \cite{bib:most-popular-mobile-frameworks-2024}, and because it allows for the use of many JavaScript frameworks, such as Angular, React, and Vue.js. \cite{bib:ionic-documentation}

\subsection{Evaluation}\label{subsec:evaluation}
The frameworks are evaluated based on the criteria established in \cref{subsec:criteria}.

\subsubsection{Flutter}\label{subsubsec:evaluation-flutter}

\paragraph{Communication Protocols}
Flutter has a well-maintained package named \emph{dio}, that can perform rest calls. \cite{bib:dio-package}
Additionally, the \emph{web\_socket\_channel} package can be used to establish WebSocket connections. \cite{bib:web-socket-channel-package}

Since this criterion is not assigned any weight, Flutter meets the requirements and passes the criterion.

\paragraph{Community}
The official Flutter GitHub repository has over 172.000 stars \cite{bib:flutter-github}, the official Flutter Discord server has over 71.000 members \cite{bib:flutter-discord}, and there are at least 523.000 results on GitHub for Flutter projects \cite{bib:flutter-github-search}.

The scores for this criterion are shown in \cref{tab:multi-criteria-analysis-flutter-community}.

\begin{htable}[
		caption = {Flutter community evaluation},
		label = {tab:multi-criteria-analysis-flutter-community},
	]{
		colspec={lllX},
	}
	Sub-criterion & Weight & Function & Normalized score \\
	Number of stars & 5 & $\bar{x} = \frac{N}{N_{\max}}$ & $\bar{x} = \frac{172000}{172000} = 1$ \\
	Number of Discord users & 5 & $\bar{x} = \frac{N}{N_{\max}}$ & $\bar{x} = \frac{71000}{308000} = 0.23$ \\
	Number of examples & 10 & $\bar{x} = \frac{N}{N_{\max}}$ & $\bar{x} = \frac{523000}{535000} = 0.98$ \\
\end{htable}

\[
	S_{\text{community}} =
	\frac{
		5 \cdot 1 +
		5 \cdot 0.23 +
		10 \cdot 0.98
	}{
		5 + 5 + 10
	} = 0.80
\]

The normalized total score for this criterion is \emph{0.80}.

\paragraph{Developer Experience}
The proficiency ratings are as follows:
\begin{itemize}
	\item \textbf{Moustafa (0.50)} - Has been working on a Flutter project at work, but feels like he does not yet have a solid understanding of the framework.
	\item \textbf{William (0.30)} - Has played around with Flutter in his free time, but does not feel like he has used the framework extensively enough to be proficient.
\end{itemize}

Debugging flutter can be done inside Android Studio using the Dart and Flutter plugins \cite{bib:flutter-debugging}.
The documentation for Flutter contains an installation guide, and covers lots of aspects of the framework, including state management and navigation \cite{bib:flutter-documentation}.

The scores for this criterion are shown in \cref{tab:multi-criteria-analysis-flutter-developer-experience}.

\begin{htable}[
		caption = {Flutter developer experience evaluation},
		label = {tab:multi-criteria-analysis-flutter-developer-experience},
	]{
		colspec={lllX},
	}
	Sub-criterion & Weight & Function & Normalized score \\
	Learning curve & 20 & $\bar{x} = \frac{x_1 + x_2}{2}$ & $\bar{x} = \frac{0.50 + 0.30}{2} = 0.40$ \\
	Debugging support & 10 & Binary (0 or 1) & 1 \\
	Documentation quality & 15 & Binary (0 or 1) & 1 \\
\end{htable}

\[
	S_{\text{developer experience}} =
	\frac{
		20 \cdot 0.40 +
		10 \cdot 1 +
		15 \cdot 1
	}{
		20 + 10 + 15
	} = 0.73
\]

The normalized total score for this criterion is \emph{0.73}.

\paragraph{Ecosystem}
There is an official Flutter plugin for JetBrains IDEs, which provides code completion and error checking \cite{bib:flutter-jetbrains-plugin}.
The Dart package repository \cite{bib:pub-dev} contains over 58.000 packages, many of which are compatible with Flutter.
For UI components, there is the built-in \emph{Material} \cite{bib:flutter-material} and \emph{Cupertino} \cite{bib:flutter-cupertino} libraries, as well as several open source libraries, such as \emph{Syncfusion} \cite{bib:flutter-syncfusion}, and \emph{GetWidget} \cite{bib:flutter-getwidget}.

The scores for this criterion are shown in \cref{tab:multi-criteria-analysis-flutter-ecosystem}.

\begin{htable}[
		caption = {Flutter ecosystem evaluation},
		label = {tab:multi-criteria-analysis-flutter-ecosystem},
	]{
		colspec={lllX},
	}
	Sub-criterion & Weight & Function & Normalized score \\
	IDE plugin support & 25 & Binary (0 or 1) & 1 \\
	Number of libraries & 15 & $\bar{x} = \frac{N}{N_{\max}}$ & $\bar{x} = \frac{58000}{58000} = 1$ \\
	Availability of UI components & 15 & Binary (0 or 1) & 1 \\
\end{htable}

\[
	S_{\text{ecosystem}} =
	\frac{
		25 \cdot 1 +
		15 \cdot 1 +
		15 \cdot 1
	}{
		25 + 15 + 15
	} = 1
\]

The normalized total score for this criterion is \emph{1.00}.

\subsubsection{React Native}\label{subsubsec:evaluation-react-native}

\paragraph{Communication Protocols}
React Native runs JavaScript, so either Fetch API \cite{bib:fetch-api} or Axios \cite{bib:axios} can be used to perform RESTful API calls.
The \texttt{WebSocket} class \cite{bib:websocket-class} from JavaScript can be used to establish WebSocket connections.

Since this criterion is not assigned any weight, React Native meets the requirements and passes the criterion.

\paragraph{Community}
The official React Native GitHub repository has over 124.000 stars \cite{bib:react-native-github}.
There is one unofficial Discord server named Reactiflux, which has over 230.000 members \cite{bib:reactiflux-discord}.
There are also three company-based Discord servers referenced on the official React Native website \cite{bib:react-native-communities}, which in total amount to 78.000 members.
Search results on GitHub for React Native projects amount to at least 535.000 repositories \cite{bib:react-native-github-search}.

The scores for this criterion are shown in \cref{tab:multi-criteria-analysis-react-native-community}.

\begin{htable}[
		caption = {React Native community evaluation},
		label = {tab:multi-criteria-analysis-react-native-community},
	]{
		colspec={lllX},
	}
	Sub-criterion & Weight & Function & Normalized score \\
	Number of stars & 5 & $\bar{x} = \frac{N}{N_{\max}}$ & $\bar{x} = \frac{124000}{172000} = 0.72$ \\
	Number of Discord users & 5 & $\bar{x} = \frac{N}{N_{\max}}$ & $\bar{x} = \frac{308000}{308000} = 1$ \\
	Number of examples & 10 & $\bar{x} = \frac{N}{N_{\max}}$ & $\bar{x} = \frac{535000}{535000} = 1$ \\
\end{htable}

\[
	S_{\text{community}} =
	\frac{
		5 \cdot 0.72 +
		5 \cdot 1 +
		10 \cdot 1
	}{
		5 + 5 + 10
	} = 0.93
\]

The normalized total score for this criterion is \emph{0.93}.

\paragraph{Developer Experience}
The proficiency ratings are as follows:

\begin{itemize}
	\item \textbf{Moustafa (0.20)} - Has never used the framework, but does have minimal experience with other JavaScript frameworks such as VueJS and Svelte.
	\item \textbf{William (0.65)} - Has built a small application using React Native, and has extensive experience with other JavaScript frameworks.
\end{itemize}

Debugging can be done inside WebStorm \cite{bib:react-native-debugging}.
The documentation for React Native contains an installation guide, and covers lots of aspects of the framework, including state management and navigation \cite{bib:react-native-documentation}.

The scores for this criterion are shown in \cref{tab:multi-criteria-analysis-react-native-developer-experience}.

\begin{htable}[
		caption = {React Native developer experience evaluation},
		label = {tab:multi-criteria-analysis-react-native-developer-experience},
	]{
		colspec={lllX},
	}
	Sub-criterion & Weight & Function & Normalized score \\
	Learning curve & 20 & $\bar{x} = \frac{x_1 + x_2}{2}$ & $\bar{x} = \frac{0.20 + 0.65}{2} = 0.43$ \\
	Debugging support & 10 & Binary (0 or 1) & 1 \\
	Documentation quality & 15 & Binary (0 or 1) & 1 \\
\end{htable}

\[
	S_{\text{developer experience}} =
	\frac{
		20 \cdot 0.43 +
		10 \cdot 1 +
		15 \cdot 1
	}{
		20 + 10 + 15
	} = 0.74
\]

The normalized total score for this criterion is \emph{0.74}.

\subsubsection{Ionic}\label{subsubsec:evaluation-ionic}

\paragraph{Communication Protocols}
Ionic, like React Native, runs JavaScript, so either Fetch API \cite{bib:fetch-api} or Axios \cite{bib:axios} can be used to perform RESTful API calls.
It can also use the \texttt{WebSocket} class \cite{bib:websocket-class} from JavaScript to establish WebSocket connections.

Since this criterion is not assigned any weight, Ionic meets the requirements and passes the criterion.

\paragraph{Community}
The official Ionic GitHub repository has over 52.000 stars \cite{bib:ionic-github}, the official Ionic Discord server has just over 11.000 members \cite{bib:ionic-discord}, and there are at least 136.000 results on GitHub for Ionic projects \cite{bib:ionic-github-search}.

The scores for this criterion are shown in \cref{tab:multi-criteria-analysis-ionic-community}.

\begin{htable}[
		caption = {Ionic community evaluation},
		label = {tab:multi-criteria-analysis-ionic-community},
	]{
		colspec={lllX},
	}
	Sub-criterion & Weight & Function & Normalized score \\
	Number of stars & 5 & $\bar{x} = \frac{N}{N_{\max}}$ & $\bar{x} = \frac{52000}{172000} = 0.30$ \\
	Number of Discord users & 5 & $\bar{x} = \frac{N}{N_{\max}}$ & $\bar{x} = \frac{11000}{308000} = 0.04$ \\
	Number of examples & 10 & $\bar{x} = \frac{N}{N_{\max}}$ & $\bar{x} = \frac{136000}{535000} = 0.25$ \\
\end{htable}

\[
	S_{\text{community}} =
	\frac{
		5 \cdot 0.30 +
		5 \cdot 0.04 +
		10 \cdot 0.25
	}{
		5 + 5 + 10
	} = 0.21
\]

The normalized total score for this criterion is \emph{0.21}.

\paragraph{Developer Experience}
The proficiency ratings are as follows:

\begin{itemize}
	\item \textbf{Moustafa (0.25)} - Has never used the framework, but does have minimal experience with other JavaScript frameworks such as VueJS and Svelte.
	      These frameworks can be used with Ionic, so some of the knowledge is transferable.
	\item \textbf{William (0.50)} - Has never used the framework, but does have experience with other JavaScript frameworks.
	      These frameworks can be used with Ionic, so some of the knowledge is transferable.
\end{itemize}

Debugging can be done inside WebStorm \cite{bib:ionic-debugging}.
The documentation for Ionic contains an installation guide, but does not cover every aspect one would expect of a framework.
It does however, link to documentation for popular JavaScript frameworks, which can be used with Ionic \cite{bib:ionic-documentation}, which each provide their own functionality like state management and navigation.

The scores for this criterion are shown in \cref{tab:multi-criteria-analysis-ionic-developer-experience}.

\begin{htable}[
		caption = {Ionic developer experience evaluation},
		label = {tab:multi-criteria-analysis-ionic-developer-experience},
	]{
		colspec={lllX},
	}
	Sub-criterion & Weight & Function & Normalized score \\
	Learning curve & 20 & $\bar{x} = \frac{x_1 + x_2}{2}$ & $\bar{x} = \frac{0.25 + 0.50}{2} = 0.38$ \\
	Debugging support & 10 & Binary (0 or 1) & 1 \\
	Documentation quality & 15 & Binary (0 or 1) & 1 \\
\end{htable}

\[
	S_{\text{developer experience}} =
	\frac{
		20 \cdot 0.38 +
		10 \cdot 1 +
		15 \cdot 1
	}{
		20 + 10 + 15
	} = 0.71
\]

The normalized total score for this criterion is \emph{0.71}.
